\documentclass[12pt, a4paper]{article}

% ── Paquetes ─────────────────────────────────────────────────────────
\usepackage[utf8]{inputenc}
\usepackage[T1]{fontenc}
% babel-spanish no disponible en esta instalación
\usepackage{amsmath, amssymb, amsthm}
\usepackage{geometry}
\usepackage{graphicx}
\usepackage{xcolor}
\usepackage{mdframed}
\usepackage{enumitem}
\usepackage{booktabs}
\usepackage{hyperref}

% ── Geometría ────────────────────────────────────────────────────────
\geometry{
    top    = 2.5cm,
    bottom = 2.5cm,
    left   = 3.0cm,
    right  = 3.0cm
}

% ── Colores ──────────────────────────────────────────────────────────
\definecolor{azulviz}{RGB}{37, 99, 235}
\definecolor{naranjviz}{RGB}{251, 146, 60}
\definecolor{verdviz}{RGB}{52, 211, 153}
\definecolor{amarillviz}{RGB}{251, 191, 36}
\definecolor{fondocaja}{RGB}{240, 245, 255}
\definecolor{bordecruz}{RGB}{37, 99, 235}
\definecolor{fondoadv}{RGB}{255, 248, 230}
\definecolor{bordeadv}{RGB}{251, 146, 60}
\definecolor{fondoteo}{RGB}{235, 250, 242}
\definecolor{bordeteo}{RGB}{52, 153, 110}

% ── Entornos ─────────────────────────────────────────────────────────
\mdfdefinestyle{vizverde}{
    backgroundcolor = fondoteo,
    linecolor       = bordeteo,
    linewidth       = 1.5pt,
    leftmargin      = 0pt,
    rightmargin     = 0pt,
    innertopmargin  = 8pt,
    innerbottommargin = 8pt,
    innerleftmargin  = 12pt,
    innerrightmargin = 12pt,
    roundcorner     = 4pt
}

\mdfdefinestyle{viznaranja}{
    backgroundcolor = fondoadv,
    linecolor       = bordeadv,
    linewidth       = 1.5pt,
    leftmargin      = 0pt,
    rightmargin     = 0pt,
    innertopmargin  = 8pt,
    innerbottommargin = 8pt,
    innerleftmargin  = 12pt,
    innerrightmargin = 12pt,
    roundcorner     = 4pt
}

\mdfdefinestyle{vizazul}{
    backgroundcolor = fondocaja,
    linecolor       = bordecruz,
    linewidth       = 1.5pt,
    leftmargin      = 0pt,
    rightmargin     = 0pt,
    innertopmargin  = 8pt,
    innerbottommargin = 8pt,
    innerleftmargin  = 12pt,
    innerrightmargin = 12pt,
    roundcorner     = 4pt
}

\newenvironment{viznota}[1]{%
    \begin{mdframed}[style=vizverde]
    \textbf{\color{bordeteo}En GradienViz \textemdash\ #1.}\quad
}{%
    \end{mdframed}
}

\newenvironment{vizadv}{%
    \begin{mdframed}[style=viznaranja]
    \textbf{\color{bordeadv}Atención:}\quad
}{%
    \end{mdframed}
}

\newenvironment{vizdef}{%
    \begin{mdframed}[style=vizazul]
}{%
    \end{mdframed}
}

% ── Comandos matemáticos ─────────────────────────────────────────────
\newcommand{\R}{\mathbb{R}}
\newcommand{\norm}[1]{\left\lVert #1 \right\rVert}
\newcommand{\grad}{\nabla}
\newcommand{\mstar}{m^{*}}
\newcommand{\bstar}{b^{*}}
\newcommand{\Jfun}{J(m,b)}

% ── Hyperref ─────────────────────────────────────────────────────────
\hypersetup{
    colorlinks = true,
    linkcolor  = azulviz,
    urlcolor   = azulviz,
    citecolor  = azulviz
}

% ── Metadatos ────────────────────────────────────────────────────────
\title{
    {\Large\bfseries Descenso por gradiente en regresión lineal}\\[0.4em]
    {\large Una introducción con Cálculo Diferencial y GradienViz}\\[0.8em]
    {\normalsize Notas de lectura para estudiantes de primer semestre}
}
\author{
    Departamento Académico de Sistemas Computacionales\\
    Universidad Autónoma de Baja California Sur
}
\date{2026}

% ──────────────────────────────────────────────────────────────────────
\begin{document}

\maketitle
\tableofcontents
\vspace{1em}

\begin{vizadv}
Este documento requiere únicamente conocimientos de \textbf{Cálculo Diferencial de una variable} (el concepto de derivada como pendiente y razón de cambio) y geometría analítica básica. Se recomienda tener abierto \textbf{GradienViz} mientras se lee para acompañar la teoría con la simulación interactiva.
\end{vizadv}

\section{Punto de partida: ajustando una recta a datos}

Cuando abres GradienViz y colocas algunos puntos en el Panel~1, la herramienta calcula automáticamente la pendiente~$\mstar$ y el intercepto~$\bstar$ de la recta que mejor se ajusta a esos datos, y marca el punto $(\mstar, \bstar)$ con una cruz amarilla~$\oplus$ en el Panel~2. Esos dos números son la solución al siguiente problema:

\begin{vizdef}
\textbf{Problema de regresión lineal simple.} Dados $n$ puntos $(x_i, y_i)$ para $i = 1, \ldots, n$, encontrar los valores $m, b \in \R$ que minimizan el error cuadrático medio:
\begin{equation}\label{eq:J}
    \Jfun = \frac{1}{n}\sum_{i=1}^{n}\bigl(m x_i + b - y_i\bigr)^2.
\end{equation}
\end{vizdef}

Para este problema existe una fórmula directa exacta. Definiendo las sumas
\[
    S_x = \sum x_i, \quad
    S_y = \sum y_i, \quad
    S_{xy} = \sum x_i y_i, \quad
    S_{x^2} = \sum x_i^2,
\]
se obtiene directamente:
\begin{equation}\label{eq:exacta}
    \mstar = \frac{n S_{xy} - S_x S_y}{n S_{x^2} - S_x^2}, \qquad
    \bstar = \frac{S_y - \mstar S_x}{n}.
\end{equation}

Si ya existe esta fórmula, \textbf{¿para qué simular un algoritmo de iteración pasito a pasito?} La respuesta es que en modelos más complejos (como las redes neuronales) no existe una fórmula directa, y la derivada es la única herramienta que nos permite ``aprender'' o encontrar la solución navegando a ciegas.

\begin{viznota}{Panel~1}
La recta gris punteada es $y = \mstar x + \bstar$, calculada analíticamente con \eqref{eq:exacta}. La recta naranja es la estimación provisional que el algoritmo va corrigiendo en cada paso hasta alcanzar a la solución exacta.
\end{viznota}

\section{La función de costo como superficie o tazón}

La expresión~\eqref{eq:J} mide el error total de la recta dados los parámetros $m$ y $b$. Si fijamos los datos $(x_i, y_i)$, cada par $(m, b)$ produce un valor de error $J$. Al graficar $J$ en función de $m$ y $b$, obtenemos una superficie tridimensional con forma de \textbf{tazón} (un paraboloide). Si expandimos la suma en~\eqref{eq:J}, observamos que para $m$ (o para $b$) la función $J$ es un polinomio de grado~2 (una parábola que abre hacia arriba). Esto garantiza que el tazón tiene \textbf{un único punto mínimo}, ubicado exactamente en $(\mstar, \bstar)$.

\begin{viznota}{Panel~2, mapa de calor}
El mapa de color viridis representa la función $\Jfun$: las zonas azul oscuro corresponden a errores pequeños (cerca del fondo del tazón) y las amarillas a errores grandes. Observa que siempre se forma un único valle sin bordes ni pozos secundarios.
\end{viznota}

\section{La derivada como guía: ¿hacia dónde movernos?}

Para encontrar el fondo del tazón desde un punto inicial cualquiera, necesitamos saber hacia dónde decrece la función. Aquí es donde entra la derivada.

La función $J(m, b)$ depende de dos variables, pero la idea es sencilla: derivamos respecto a una sola variable a la vez, \textbf{tratando la otra como si fuera una constante}. Esta operación\textemdash que practicarás formalmente al final del semestre bajo el nombre de \emph{derivada parcial}\textemdash se aplica igual que cualquier derivada de polinomio.

\medskip
\textbf{Derivada respecto a $m$ (tratando $b$ como constante):}
\begin{equation}\label{eq:gm}
    \frac{\partial J}{\partial m} = \frac{2}{n}\sum_{i=1}^{n}(m x_i + b - y_i)\,x_i.
\end{equation}

\textbf{Derivada respecto a $b$ (tratando $m$ como constante):}
\begin{equation}\label{eq:gb}
    \frac{\partial J}{\partial b} = \frac{2}{n}\sum_{i=1}^{n}(m x_i + b - y_i).
\end{equation}

El símbolo $\partial$ (``d redonda'') es la notación estándar para estas derivadas parciales; por ahora basta saber que se calculan exactamente igual que una derivada ordinaria, fijando la otra variable.

\subsection{El gradiente}

El \textbf{gradiente}, denotado por $\grad J(m, b)$, es simplemente un vector que agrupa ambas derivadas parciales:
\[
    \grad J(m, b) =
    \begin{pmatrix} \partial J/\partial m \\[4pt] \partial J/\partial b \end{pmatrix}
    = \frac{2}{n}\begin{pmatrix}
        \displaystyle\sum_{i=1}^{n}(m x_i + b - y_i)\,x_i \\[8pt]
        \displaystyle\sum_{i=1}^{n}(m x_i + b - y_i)
    \end{pmatrix}.
\]
Geométricamente:
\begin{enumerate}[label=(\roman*)]
    \item La derivada nos da la pendiente o inclinación de la superficie en ese punto.
    \item El vector gradiente $\grad J$ apunta en la dirección de \textbf{máximo crecimiento} del error.
    \item Por lo tanto, el sentido opuesto, $-\grad J$, apunta directamente en la dirección de \textbf{mayor descenso} (hacia el fondo del tazón).
\end{enumerate}

\begin{viznota}{Panel~2, campo vectorial}
Las flechas pequeñas en el Panel~2 representan el vector de descenso $-\grad J$ en cada punto. Nota cómo todas las flechas apuntan hacia la cruz amarilla~$\oplus$ y se hacen más cortas cerca del mínimo, reflejando que las derivadas se vuelven cero en el punto óptimo.
\end{viznota}

\section{El algoritmo de descenso por gradiente}

La idea del descenso por gradiente es simple: comenzar en una estimación inicial $(m_0, b_0)$ y dar pasos en la dirección opuesta a la derivada.

\begin{vizdef}
\textbf{Regla de actualización.} En cada paso $t$, corregimos los parámetros
restando una fracción de la derivada parcial correspondiente,
evaluada en los valores actuales $(m^{(t)}, b^{(t)})$:
\begin{align*}
    m^{(t+1)} &= m^{(t)}
      - \frac{2\eta}{n}\sum_{i=1}^{n}\bigl(m^{(t)} x_i + b^{(t)} - y_i\bigr)\,x_i, \\[6pt]
    b^{(t+1)} &= b^{(t)}
      - \frac{2\eta}{n}\sum_{i=1}^{n}\bigl(m^{(t)} x_i + b^{(t)} - y_i\bigr).
\end{align*}
O en forma compacta, usando la notación de gradiente:
\[
    \begin{pmatrix} m^{(t+1)} \\ b^{(t+1)} \end{pmatrix}
    =
    \begin{pmatrix} m^{(t)} \\ b^{(t)} \end{pmatrix}
    - \eta\,\grad J\!\left(m^{(t)}, b^{(t)}\right).
\]
La constante $\eta > 0$ (\textit{eta}) se conoce como \textbf{tasa de aprendizaje}
(\textit{learning rate}).
\end{vizdef}

\subsection{El rol de la tasa de aprendizaje \texorpdfstring{$\eta$}{eta}}

La tasa de aprendizaje $\eta$ determina el tamaño del paso en cada iteración:

\begin{center}
\begin{tabular}{lll}
    \toprule
    Régimen              & Comportamiento                        & Causa \\
    \midrule
    $\eta$ muy pequeño  & Convergencia muy lenta                & Pasos demasiado cortos \\
    $\eta$ adecuado     & Convergencia suave hacia $(\mstar,\bstar)$ & Paso equilibrado \\
    $\eta$ demasiado grande & Oscilaciones o divergencia          & El paso ``rebasa'' el valle \\
    \bottomrule
\end{tabular}
\end{center}

\begin{viznota}{slider $\eta$}
Modifica el deslizador de $\eta$. Si asignas un valor muy pequeño el algoritmo tardará cientos de iteraciones en llegar a la solución; si pruebas un valor alto (p. ej., $\eta > 0.3$), verás cómo las trayectorias rebotan y se disparan fuera de la pantalla.
\end{viznota}

\section{Geometría del error y el parámetro \texorpdfstring{$\kappa$}{kappa}}

Si el tazón del error fuera perfectamente circular, las derivadas apuntarían directo al centro y el algoritmo convergiría en pocas iteraciones. Sin embargo, dependiendo de la distribución de los puntos en el Panel~1, el tazón puede deformarse hasta convertirse en un valle muy estrecho y alargado.

GradienViz calcula un indicador de esa deformación llamado \textbf{número de condición} ($\kappa$), que mide la razón entre la curvatura más pronunciada y la más suave del tazón:

\begin{itemize}
    \item \textbf{$\kappa \approx 1$\ (tazón redondo):}
    Ambas curvaturas son similares. El gradiente apunta casi recto al mínimo
    y el algoritmo converge en pocas iteraciones.

    \item \textbf{$\kappa = 100$\ (valle moderado):}
    El tazón es 100 veces más ancho en una dirección que en la otra.
    El gradiente toma un camino oblicuo, necesitando más iteraciones.

    \item \textbf{$\kappa \gg 1000$\ (valle estrecho):}
    La pendiente lateral del valle es miles de veces mayor que la del fondo.
    El algoritmo avanza en zigzag de pared en pared, con convergencia muy lenta
    o nula dentro del límite de iteraciones.
\end{itemize}

\begin{viznota}{Panel 1, esquina superior derecha}
Arrastra los puntos en el Panel 1 para alinearlos casi de forma vertical (mismo valor de $x$). Verás cómo el valor de $\kappa$ se dispara y se torna naranja. Geométricamente, si todas las $x$ son iguales, es imposible determinar la pendiente $m$ de una recta, por lo que el algoritmo sufre para encontrar una solución.
\end{viznota}

\section{De la recta a las redes neuronales}

Para una recta, la fórmula exacta~\eqref{eq:exacta} es instantánea y no requiere iteraciones. ¿Por qué el descenso por gradiente es entonces el pilar fundamental de la Inteligencia Artificial?

\subsection{De 2 parámetros a miles de millones}

En regresión lineal simple solo ajustamos 2 parámetros ($m$ y $b$). Una red neuronal para clasificar imágenes o un modelo de lenguaje (como GPT) puede tener entre miles y cientos de \textbf{miles de millones} de parámetros. A esa escala no existen fórmulas analíticas cerradas, y el descenso por gradiente es la única alternativa computable.

\subsection{La regla de aprendizaje es idéntica}

Aunque una red neuronal combine múltiples capas y funciones no lineales, la regla con la que aprende parámetro por parámetro en cada paso es exactamente la misma que acabas de ver:
\[
    \text{Parámetro nuevo} = \text{Parámetro actual}
    - \eta \times \left(\frac{\partial\,\text{Error}}{\partial\,\text{Parámetro}}\right).
\]
En las redes neuronales hay miles de parámetros y la función de error pasa por muchas capas de operaciones anidadas. Para calcular cada derivada parcial se aplica la \textbf{regla de la cadena}\textemdash la misma que usas en Cálculo~I para derivar composiciones de funciones\textemdash de manera sistemática desde la salida de la red hasta cada uno de sus parámetros. Este procedimiento se llama \textbf{retropropagación} (\textit{backpropagation}) y es el corazón computacional de todo el aprendizaje profundo.

\subsection{Resumen comparativo}

\begin{center}
\small
\begin{tabular}{lll}
    \toprule
    & \textbf{Regresión lineal (GradienViz)} & \textbf{Red neuronal} \\
    \midrule
    Parámetros a ajustar      & 2 ($m$ y $b$)               & Miles a miles de millones \\
    ¿Se puede visualizar?     & Sí (Panel~2, 2D)            & No (demasiadas dimensiones) \\
    Forma del error           & Tazón suave (único mínimo)  & Superficie con múltiples valles \\
    ¿Existe fórmula exacta?   & Sí \eqref{eq:exacta}       & No \\
    Método de ajuste          & Descenso por gradiente      & Descenso por gradiente \\
    Herramienta clave         & Derivadas parciales         & Derivadas parciales + regla de la cadena \\
    \bottomrule
\end{tabular}
\end{center}

\bigskip
Cuando presionas \textbf{Train} en GradienViz, estás observando en dos dimensiones exactamente el mismo principio con el que se entrenan los modelos de inteligencia artificial más avanzados del mundo: ajustar parámetros paso a paso siguiendo la dirección que señalan las derivadas.

\end{document}
